\lysh{38829}{4}

\begin{alist}
\item Η Μαρία διαβάζει κάθε μέρα $5$ σελίδες περισσότερες από την προηγούμενη. Έτσι ο αριθμός των σελίδων που θα διαβάσει τη $\nu$-οστή μέρα ακολουθεί μία αριθμητική πρόοδο με πρώτο όρο $\alpha_1 = 20$ και διαφορά $\omega = 5$ η οποία δίνεται από τον τύπο $\alpha_\nu = \alpha_1 + (\nu - 1)\omega$. Άρα ο αριθμός των σελίδων που θα διαβάσει την 5η μέρα είναι:
\[\alpha_5 = 20 + (5 - 1) \cdot 5 = 20 + 4 \cdot 5 = 20 + 20 = 40\]

\item Για να δούμε αν θα προλάβει να διαβάσει η Μαρία την ύλη του Μ1 μαθήματος, που είναι $296$ σελίδες, θα πρέπει να βρούμε πόσες σελίδες θα έχει διαβάσει η Μαρία μέχρι και την 8η μέρα. Θα υπολογίσουμε το άθροισμα των σελίδων από τον τύπο
\[S_\nu = \dfrac{\nu}{2}[2\alpha_1 + (\nu - 1) \cdot \omega].\]
Άρα $S_8 = \dfrac{8}{2}[2 \cdot 20 + (8 - 1) \cdot 5] = 4 \cdot (40 + 35) = 4 \cdot 75 = 300$ σελίδες.

Άρα, σε $8$ μέρες αφότου θα ξεκινήσει το διάβασμα θα μπορέσει να ολοκληρώσει την ύλη του πρώτου μαθήματος, που είναι $296$ σελίδες, αφού $296 < 300$.

\item Όλη η ύλη είναι $296 + 274 + 400 = 970$ σελίδες. Για να βρούμε πόσες μέρες η Μαρία θα χρειαστεί συνολικά για να διαβάσει όλη την ύλη των $970$ σελίδων, θα χρησιμοποιήσουμε και πάλι τον τύπο $S_\nu$, για να βρούμε όμως το $\nu$.
\begin{align*}
S_\nu = 970 \Leftrightarrow 970 = \dfrac{\nu}{2}[2 \cdot 20 + (\nu - 1) \cdot 5] \Leftrightarrow 1940 = \nu[40 + (\nu - 1) \cdot 5] \ &\Leftrightarrow \\
1940 &= 40\nu + 5\nu^2 - 5\nu \Leftrightarrow 5\nu^2 + 35\nu - 1940 = 0 \Leftrightarrow \nu^2 + 7\nu - 388 = 0 \\
\Delta &= 49 + 1552 = 1601 \\
\nu_{1,2} &= \dfrac{-7 \pm \sqrt{1601}}{2} \simeq \dfrac{-7 \pm 40}{2}\text{, οπότε}
\end{align*}
$\nu_1 = \dfrac{33}{2}$, $\nu_2 = \dfrac{-47}{2}$, δηλαδή
$\nu_1 \simeq 16,5$, $\nu_2 \simeq -23,5$.
Το $\nu_2 \simeq -23,5$ απορρίπτεται. Άρα η Μαρία θα χρειαστεί $17$ ημέρες για να διαβάσει την ύλη και των $3$ μαθημάτων ή των $970$ σελίδων.

\item Τα $2$ πρώτα μαθήματα Μ1 και Μ2 έχουν συνολικά $274 + 296 = 570$ σελίδες.
Αρχικά θα βρούμε σε πόσες μέρες θα ολοκληρώσει η Μαρία τις $570$ σελίδες υπολογίζοντας το $\nu$ στην εξίσωση $S_\nu = 570$.
\begin{align*}
S_\nu = 570 \Leftrightarrow \dfrac{\nu}{2}[2 \cdot 20 + (\nu - 1) \cdot 5] = 570 \Leftrightarrow \nu \cdot [40 + (\nu - 1) \cdot 5] = 1140 \ &\Leftrightarrow \\
40\nu + 5\nu^2 - 5\nu &= 1140 \Leftrightarrow 5\nu^2 + 35\nu - 1140 = 0 \Leftrightarrow \nu^2 + 7\nu - 228 = 0 \\
\Delta &= 961 \\
\nu_{1,2} &= \dfrac{-7 \pm \sqrt{961}}{2} = \dfrac{-7 \pm 31}{2}
\end{align*}
$\nu_1 = \dfrac{-7 + 31}{2} = \dfrac{24}{2} = 12$ ή $\nu_2 = \dfrac{-7 - 31}{2} = \dfrac{-38}{2} = -19$, που απορρίπτεται.

Άρα για να ολοκληρώσει την ύλη των $2$ πρώτων μαθημάτων Μ1 και Μ2, δηλαδή τις $570$ πρώτες σελίδες, θα χρειαστεί $12$ ημέρες. Από το προηγούμενο ερώτημα είδαμε ότι χρειάζεται $17$ ημέρες για να διαβάσει την ύλη και των τριών μαθημάτων συνολικά, άρα, αν από τις $17$ ημέρες αφαιρέσουμε τις $12$ που θα χρειαστεί για τα Μ1 και Μ2, τότε για το Μ3 θα χρειαστεί συνολικά $5$ ημέρες, και συγκεκριμένα τις τελευταίες $5$ ημέρες του διαβάσματός της, ακολουθώντας πάντα τον ίδιο ρυθμό.
Εναλλακτικά, θα μπορούσαμε να υπολογίσουμε το άθροισμα:
$\alpha_{13} + \alpha_{14} + \alpha_{15} + \alpha_{16} = 80 + 85 + 90 + 95 = 350$ και μένουν και άλλες $50$ σελίδες, για να τις διαβάσει την 5η ημέρα στη σειρά από αυτές.
Άρα θα χρειαστεί $5$ ημέρες για να διαβάσει το τελευταίο μάθημα, και την 5η ημέρα από αυτές θα της μένουν να διαβάσει μόνο $50$ σελίδες αντί για $100$.
Άρα η ημερομηνία έναρξης διαβάσματος του Μ3 θα πρέπει να είναι το πολύ στις 9 Ιουνίου, ώστε να έχει τελειώσει τουλάχιστον $2$ μέρες πριν από τις 15/6, που είναι η έναρξη των εξετάσεων.
\end{alist}
